%==============================================================================
\section{Introduction}
\label{sec:intro}
%==============================================================================

\IEEEPARstart{L}{arge} language models (LLMs) are increasingly asked to reason over contexts that do not fit in, or must not leave, a single machine: long documents, multi-party logs, or private records held by different edge devices. Context-parallel decoding addresses this by sharding the key--value (KV) cache along the sequence dimension: every shard computes a partial attention for the current query and the partial results are merged by a log-sum-exp (LSE) reduction~\cite{lin2024infinite,acharya2025star,bhatia2025helix,shyam2024tree}. The merge message of a shard is a single vector and a single scalar~\cite{acharya2025star}, and the merge itself is a \emph{normalized aggregation}
\begin{equation}
\vo=\frac{\sum_{j}\vn_j}{\sum_{j}d_j}=\frac{\sum_j d_j\obar_j}{\sum_j d_j},\qquad d_j=\sum_{i\in\mathcal{S}_j}e^{s_i-c},
\label{eq:intro_merge}
\end{equation}
where $s_i$ are the attention logits of the tokens in shard $\mathcal{S}_j$, $c$ is a common shift and $\obar_j$ is the shard's local attention output. The same structure appears in kernel (Nadaraya--Watson) regression of spatial fields, where $d_j$ is a kernel of the distance between the query location and sensor~$j$, and in product-of-experts Gaussian-process regression, where $d_j$ is the local posterior precision~\cite{sato2022gpr,sato2026awa}.

Because both the numerator and the denominator of~\eqref{eq:intro_merge} are sums, over-the-air computation (AirComp)~\cite{nazer2007computation,goldenbaum2015nomographic,sahin2023survey} can produce them in one shot by letting all devices transmit simultaneously, after which the fusion center (FC) divides. For attention, one aggregation instance per layer and head occupies $d_v+1$ real channel uses regardless of the number of devices~$N$. The difficulty is that the weights are \emph{private}: $d_j$ is computed at device~$j$ from the broadcast query and its own data, it is unknown to the FC, it is heavy-tailed, and it changes with every query. Weight-aware AirComp designs~\cite{cao2020power,liu2020scaling,yang2026airmoe,zhang2026logit} assume that the transmit weights are known to the server, and the only ratio-type AirComp scheme with locally computed weights~\cite{sato2026awa} collects the message norms digitally before every slot and clips the weights with bounds tuned by Bayesian optimization, without a closed-form analysis. With private weights, the scale that every device must share is limited by the \emph{mass-weighted weakest link} $\max_j d_j^2/g_j$, where $g_j$ is the effective channel gain---a quantity the FC cannot observe.

This paper develops the theory and the designs for this setting and validates them on two applications: sharded-KV LLM decoding and spatial-field estimation.

\textbf{Contributions.}
\begin{itemize}
\item \emph{Model.} We formalize private-weight normalized AirComp under the per-device peak-power convention of~\cite{yang2026airmoe}, with an $M$-antenna FC whose receive combiner depends on the channels only, and an anchor block held by the FC (Section~\ref{sec:model}). The model covers the LSE merge of context-parallel decoding and kernel/product-of-experts regression, and it treats \emph{what the FC learns} as a design dimension.
\item \emph{Extreme-value analysis.} We show that the mass-weighted inversion bottleneck is exactly Fr\'echet with scale $S=\sum_j d_j^2/\beta_j$ under Rayleigh fading and heterogeneous path loss, and give the product law under Nakagami-$m$ fading (Section~\ref{sec:analysis}). For homogeneous links $S=D^2/\Neff$, so the participation ratio $\Neff$ of the weights acts as an SNR multiplier; the law reduces to~\cite[Cor.~2]{liu2020scaling} and~\cite[Lemma~5]{chen2024view} for equal weights. We derive the denominator outage in closed form, the conditions under which the ratio estimator has finite moments, and the saturation law of any weight-blind scalar design, whose \emph{price of blindness} grows with the dispersion of $S$ across queries; on attention weights the best blind scale needs 19--28~dB.
\item \emph{Converse and designs.} An anonymous ladder handshake lets the FC learn only a quantized $\max_j d_j^2/g_j$. We prove that, among all scales that never saturate and reveal at most $k$ bits per instance, the ladder is optimal and loses $W2^{-k-1}$~dB on a range of $W$~dB, which yields a privacy--accuracy trade-off (Section~\ref{sec:design}). Known-weight designs such as~\cite{yang2026airmoe} do better because they let the bottleneck devices saturate, which a ratio with an exact denominator cannot tolerate. We remove this obstacle with two components. In the \emph{saturation-aware split ratio}, the numerator is sent over the air and the devices securely aggregate~\cite{bonawitz2017secagg} the weights with which they actually arrive, so that the ratio remains a convex combination under any saturation. With the \emph{noise-matched level report}, every device announces the smallest noise level at which its own saturation distortion stays below $\kappa$ noise variances; the resulting bottleneck has a closed-form law, is bounded instead of Fr\'echet, and needs no weight at the FC. The FC learns one quantized maximum and one sum per instance.
\item \emph{Validation.} We verify every closed form by Monte Carlo. On sharded-KV decoding of three LLMs with up to 128 devices and 32k-token contexts, the proposed design needs 5--12~dB less SNR than a non-saturating split ratio and 2--9~dB less than known-weight aggregation at the genie scale; it is on par with or better than the threshold design of AirMoE~\cite{yang2026airmoe} and stays within 0.2--4.1~dB of that design strengthened with our saturation-aware denominator. We report paired confidence intervals, a design study of both components, one-round alternatives, channel impairments, task-level retrieval, the ring size that a secure digital merge needs at the same accuracy, and a latency ledger built on it. On distributed regression of a spatial field with up to 1024 sensors, the design outperforms adaptively weighted averaging~\cite{sato2026awa} without tuned clipping or truncation (Sections~\ref{sec:regime} and~\ref{sec:exp}).
\end{itemize}

\textbf{Related work.} Power control for AirComp in fading channels is well understood when the data to be summed are normalized: threshold-based optimal policies and scaling laws~\cite{cao2020power,liu2020scaling}, multi-antenna receivers~\cite{zhu2019mimo,chen2024view}, and imperfect channel state information (CSI)~\cite{evgenidis2024imperfect}. Ratio-type AirComp with an over-the-air denominator was introduced for product-of-experts regression~\cite{sato2022gpr} and adaptively weighted averaging~\cite{sato2026awa}. AirComp has recently been applied to LLM inference: tensor-parallel all-reduce~\cite{zhang2025jstsp}, mixture-of-experts aggregation with optimal threshold power control~\cite{yang2026airmoe}, and logit-aware MIMO transceivers~\cite{zhang2026logit}; in all of them the aggregation weights are computed and known at the server, and none estimates a denominator. The threshold scale of~\cite{yang2026airmoe} is our strongest baseline, and the noise-matched report can be read as its decentralized counterpart for weights that the server must not see. Systems work on context-parallel decoding~\cite{lin2024infinite,acharya2025star,bhatia2025helix,shyam2024tree,meta2025cp} merges digitally and motivates sharding by memory capacity and data placement, as we do. To our knowledge, the LSE merge of attention has not been studied as an over-the-air ratio with private weights.
